\chapter{Desarrollo}
\section{Analizador léxico}

% ############################################################
% #                                                          #
% #               ALFABETO DE ENTRADA                        #
% #                                                          #
% ############################################################

\subsection{Alfabeto de entrada}
\begin{itemize}[leftmargin=*]
  \item Letras minúsculas: abcdefghijklmnñopqrstuvwxyz
  \item Letras mayúsculas: ABCDEFGHIJKLMNÑOPQRSTUVWXYZ
  \item Números: 0123456789
  \item Símbolos especiales: \verb~+-*/=<>!&|;:,.$#%@^%¿?¡!(){}[]_~
  \item A literales str se agrega: áéíóúÁÉÍÓÚ
\end{itemize}

% ############################################################
% #                                                          #
% #               TOKENS RECONOCIDOS                         #
% #                                                          #
% ############################################################

\subsection{Tokens reconocidos}

% ============================================================
%                IDENTIFICADOR MET_VAR
% ============================================================

\subsubsection{Identificador de método o variable}
\begin{itemize}[leftmargin=*]
  \item Expresión regular: \verb~[a-z] . ([a-z] | [A-Z] | [0-9] | [_])*~
\end{itemize}

% ============================================================
%                IDENTIFICADOR CLASE
% ============================================================

\subsubsection{Identificador de clase}
\begin{itemize}[leftmargin=*]
  \item Expresión regular: \verb~[A-Z] . ([a-z] | [A-Z] | [0-9] | [_])*~
\end{itemize}

% ============================================================
%                PALABRAS RESERVADAS
% ============================================================

\subsubsection{Palabras reservadas}
\begin{itemize}[leftmargin=*]
  \item Expresiones regulares:
  \subitem \verb~class: c.l.a.s.s~
  \subitem \verb~if: i.f~
  \subitem \verb~impl: i.m.p.l~
  \subitem \verb~in: i.n~
  \subitem \verb~else: e.l.s.e~
  \subitem \verb~false: f.a.l.s.e~
  \subitem \verb~fn: f.n~
  \subitem \verb~for: f.o.r~
  \subitem \verb~ret: r.e.t~
  \subitem \verb~while: w.h.i.l.e~
  \subitem \verb~ret: r.e.t~
  \subitem \verb~true: t.r.u.e~
  \subitem \verb~new: n.e.w~
  \subitem \verb~nil: n.i.l~
  \subitem \verb~st: s.t.~
  \subitem \verb~self: s.e.l.f~
  \subitem \verb~pub: p.u.b~
  \subitem \verb~void: v.o.i.d~
\end{itemize}

% ============================================================
%                LITERALES ENTEROS
% ============================================================

\subsubsection{Literales enteros}
\begin{itemize}[leftmargin=*]
  \item Expresión regular: \verb~[0-9]+~
\end{itemize}

% ============================================================
%                LITERALES STRING
% ============================================================

\subsubsection{Literales string}
\shorthandoff{"}

\begin{itemize}[leftmargin=*]
    \item Expresión regular para literales \texttt{Str}:
    \[
    \text{"} \cdot
    \left(
        (\Sigma - \{\text{"}\})
        \cup
        \{á,é,í,ó,ú,Á,É,Í,Ó,Ú\}
        \cup
        \left(
            \backslash \cdot
            (\text{"} \mid \text{n} \mid \text{t}
            \mid \backslash \mid \text{r})
        \right)
    \right)^*
    \cdot \text{"}
    \]

\end{itemize}

\shorthandon{"}

% ============================================================
%                ESPACIOS EN BLANCO
% ============================================================

\subsubsection{Espacios en blanco}
\begin{itemize}[leftmargin=*]
  \item Expresión regular: \verb~(" " | \n | \t | \r | \v)*~
\end{itemize}

% ============================================================
%                COMENTARIOS
% ============================================================

\subsubsection{Comentarios}
\begin{itemize}[leftmargin=*]
  \item Expresión regular para comentario simple:
  \texttt{//} $(\Sigma - \{\backslash n\})^* \backslash n$

  \item Expresión regular para comentario multilínea:
  \texttt{/*} $\Sigma^*$ \texttt{*/}
\end{itemize}

% ============================================================
%                CLASES PREDEFINIDAS DE TINYS
% ============================================================

\subsubsection{Clases predefinidas}
\begin{itemize}[leftmargin=*]
  \item Expresiones regulares:
  \subitem \verb~Object: O.b.j.e.c.t~
  \subitem \verb~IO: I.O~
  \subitem \verb~Iterator: I.t.e.r.a.t.o.r~
  \subitem \verb~Array: A.r.r.a.y~
  \subitem \verb~Int: I.n.t~
  \subitem \verb~Str: S.t.r~
  \subitem \verb~Bool: B.o.o.l~
\end{itemize}

% ============================================================
%                OPERADORES DE TINYS
% ============================================================

\subsubsection{Operadores}
\begin{itemize}[leftmargin=*]
  \item Expresiones regulares:

  \subitem \textbf{Operadores aritméticos:}
  \subitem \verb~Suma: +~
  \subitem \verb~Resta: -~
  \subitem \verb~Multiplicación: *~
  \subitem \verb~División: /~
  \subitem \verb~Incremento: +.+~
  \subitem \verb~Decremento: -.-~

  \subitem \textbf{Operadores relacionales:}
  \subitem \verb~Menor: <~
  \subitem \verb~Menor o igual: <.=~
  \subitem \verb~Mayor: >~
  \subitem \verb~Mayor o igual: >.=~
  \subitem \verb~Igualdad: =.=~
  \subitem \verb~Desigualdad: !.=~

  \subitem \textbf{Operadores lógicos:}
  \subitem \verb~AND: &.&~
  \subitem \verb~OR: |.|~
  \subitem \verb~NOT: !~

  \subitem \textbf{Operador de asignación:}
  \subitem \verb~Asignación: =~
\end{itemize}

% ============================================================
%                SÍMBOLOS DE TINYS
% ============================================================

\subsubsection{Símbolos}
\begin{itemize}[leftmargin=*]
  \item Expresiones regulares:
  \subitem \verb~Paréntesis izquierdo: (~
  \subitem \verb~Paréntesis derecho: )~
  \subitem \verb~Llave izquierda: {~
  \subitem \verb~Llave derecha: }~
  \subitem \verb~Corchete izquierdo: [~
  \subitem \verb~Corchete derecho: ]~
  \subitem \verb~Punto y coma: ;~
  \subitem \verb~Coma: ,~
  \subitem \verb~Punto: .~
  \subitem \verb~Dos puntos: :~
\end{itemize}